%% maestro2tex -- generated figure; do not edit by hand.
\providecommand{\MaestroRoot}{include/analog/}
\begin{figure}[htbp]
  \centering
  {\pgfplotsset{compat=1.18}%
  \begin{tikzpicture}
    \begin{axis}[
      width=0.82\linewidth, height=6.2cm,
      xlabel={Test-source current $i_{\mathrm{we}}$ drawn out of the WE node~[\si{\micro\ampere}]}, ylabel={Output voltage $V_{\mathrm{OUT}}$~[\si{\volt}]},
      grid=both,
      major grid style={line width=0.2pt, draw=black!14},
      minor grid style={line width=0.2pt, draw=black!7},
      minor tick num=1,
      axis line style={line width=0.4pt, draw=black!45},
      tick style={line width=0.4pt, draw=black!45},
      tick align=outside,
      label style={font=\small}, tick label style={font=\small},
      enlarge x limits=0.02, enlarge y limits=0.10,
      every axis plot/.append style={line join=round, no marks},
      legend style={font=\footnotesize, draw=none, fill=none,
                    at={(0.5,1.02)}, anchor=south, legend columns=-1,
                    /tikz/every even column/.append style={column sep=9pt}},
      legend cell align=left,
      scaled y ticks=false, scaled x ticks=false,
      y tick label style={/pgf/number format/fixed,
                          /pgf/number format/precision=1},
      ymin=0, ymax=2.5
    ]
      \addplot[mtxA, line width=1.1pt, solid] table {\MaestroRoot data/tia_dc_00.dat};
      \addlegendentry{tt}
      \addplot[mtxB, line width=1.1pt, dashed] table {\MaestroRoot data/tia_dc_01.dat};
      \addlegendentry{ff}
      \addplot[mtxC, line width=1.1pt, dotted] table {\MaestroRoot data/tia_dc_02.dat};
      \addlegendentry{fs}
      \addplot[mtxD, line width=1.1pt, dashdotted] table {\MaestroRoot data/tia_dc_03.dat};
      \addlegendentry{sf}
      \addplot[mtxE, line width=1.1pt, densely dashed] table {\MaestroRoot data/tia_dc_04.dat};
      \addlegendentry{ss}
    \end{axis}
  \end{tikzpicture}}
  \caption[{Output voltage against the swept test-source current, over the five process corners.}]{Output voltage against the swept test-source current, over the five process corners. Conditions are $R_f = \SI{560}{\kilo\ohm}$ with the working electrode held at $v_{\mathrm{cm}} = \SI{1.25}{\volt}$. The bench source draws its current \emph{out} of the WE node, so $i_{\mathrm{we}} = -I_{\mathrm{WE}}$ against the convention of Equation~(\ref{eq:tia-transfer}) and the plotted slope is $+R_f$: the amplifier sources this current through $R_f$ and the output rises. The \SI{\pm 2.26}{\micro\ampere} sweep takes the output to within \SI{12}{\milli\volt} of \texttt{vdd} and \SI{1.2}{\milli\volt} of \texttt{vss}. The corners overlay to the line width because the slope is the feedback resistor, an ideal component in this bench.}
  \label{fig:tia-dc}
\end{figure}
